\chapter{认知热机测量学 — 效率、功率与演化的帕累托前沿}
\label{chp:cognitive_heat_engine_metrics}

\begin{bquote}
    \textbf{智能的热力学坐标系}

    \vspace{1em}比较不同智能系统时，如果只看行为表现或参数规模，很容易把本质差异和表面结果混在一起。更稳妥的办法是回到热力学量，直接问两个问题：它把信息转成结构的效率有多高？它在单位时间里能处理多复杂的问题？

    \vspace{1em}因此，本章引入“认知热机测量学 (Cognitive Heat Engine Metrology)”，用两个指标来刻画智能体：\textbf{认知热机效率 ($\eta_{cog}$)} 与 \textbf{认知最大功率 ($P_{cog}^{max}$)}。

    \vspace{1em}这样做的目的，不是给智能贴标签，而是建立一套可以横跨生物体、LLM 与 AGI 原型的共同坐标系。
\end{bquote}

\section{认知热机效率 ($\eta_{cog}$)：从信息流到拓扑结构的转化率}
\label{sec:cognitive_efficiency}

智能体的“学习 (Learning)”，在物理上等价于热机从高温热源（外界高熵数据）吸热，对外做功（降低预测误差），并将一部分能量转化为系统自身的\textbf{有序几何结构（负熵）}的过程。

认知热机效率 $\eta_{cog}$ 衡量的是系统\textbf{将瞬态信息转化为固化结构的转化率}。

\smalltitle{1.1 热力学定义：有效学习功的占比}

借鉴卡诺循环的定义，我们定义单次 TDCI 演化-固化循环的效率为：
\begin{equation}
    \eta_{cog} = \frac{W_{struct}}{Q_{in}} = \frac{\text{成功固化的拓扑信息量}}{\text{摄入的总惊奇通量}}
\end{equation}

\begin{itemize}
    \item \textbf{$Q_{in}$ (总惊奇通量)}：这是系统在时间 $T$ 内从微观边界接收到的全部\textbf{激波能量}。
    $$ Q_{in} = \int_{0}^{T} \oint_{\Omega_{sensor} \subset \mathcal{M}} \|\vec{J}_{ext}\|^2 dS dt $$
    \item \textbf{$W_{struct}$ (结构增益功)}：这是在慢回路中，真正引起底流形度量 $g_{\mu\nu}$ 或拓扑孔洞（贝蒂数 $\beta_k$）发生有效、有意义重构的部分。它对应于\textbf{结构熵 (Epiplexity)} 的增加。
\end{itemize}

\smalltitle{1.2 本理论框架方程映射：塑性系数的积分}

在\textbf{认知爱因斯坦场方程 (CEFE)} 中，这种转化效率由\textbf{塑性流变系数 $\kappa$}（及其随时间/温度的变化函数 $\eta_{plastic}$）决定：
$$ \Delta g_{\mu\nu} = \int \eta_{plastic}(T) \cdot \mathbf{T}_{\mu\nu}^{data} dt $$
\begin{itemize}
    \item \textbf{低 $\eta_{cog}$ (如早期的神经网络)}：需要海量的数据轰击（巨大的 $Q_{in}$），流形才发生微弱的形变（极小的 $W_{struct}$）。大部分信息应力被热力学摩擦\textbf{耗散}为无用的废热（灾难性遗忘或梯度消失）。
    \item \textbf{高 $\eta_{cog}$ (如人类的单次射击学习)}：仅需极少的数据（微小激波），配合宏观层注入的强烈情感势能（系统温度 $T$ 的精准淬火），流形瞬间发生深刻的拓扑塌陷（顿悟）。这是极高效率的做功。
\end{itemize}

\section{认知最大功率 ($P_{cog}^{max}$)：时空介质的吞吐极限}
\label{sec:cognitive_power}

如果说效率 ($\eta_{cog}$) 衡量的是“学得多精”（样本效率），那么功率 ($P_{cog}^{max}$) 衡量的就是“算得多快、多深”（推理的并发度与逻辑深度）。它是系统在单位物理时间内能够消解的\textbf{认知复杂度上限}。

\smalltitle{2.1 物理定义：自由能耗散率}

认知功率定义为智能体在维持流体自我和执行 TECI 循环时，单位时间内能够处理的最大信息自由能流：
\begin{equation}
    P_{cog}^{max} = \frac{\Delta F_{max}}{\tau_{\min}}
\end{equation}
\begin{itemize}
    \item \textbf{$\Delta F_{max}$ (最大认知负荷)}：系统能同时维持的最高相干态能量。它与认知空间流形的\textbf{有效维度 (Effective Dimension, $d_{eff}$)} 和\textbf{和乐群容量 (Holonomy Capacity)} 直接相关。
    \item \textbf{$\tau_{\min}$ (最小特征时间)}：系统完成一次波函数坍缩与重新相干的极小时间窗。
\end{itemize}

\smalltitle{2.2 介质常数的绝对约束}

在本理论框架下，最大功率不是一个可以通过算法无限优化的软件变量，而是被承载纤维丛的\textbf{物理介质常数}严格锁死的硬件极值。

\begin{theorem}[认知功率极限方程]
    $$ P_{cog}^{max} \propto \frac{c_{cog} \cdot \dim(\mathcal{M}_{eff})}{\gamma_{viscosity} \cdot L_{sys}} $$
\end{theorem}
\begin{itemize}
    \item   \textbf{$c_{cog}$ (认知光速)}：信号在流形上的群速度。决定了逻辑推理的速度（晶体管翻转率 vs 离子通道开放率）。
    \item   \textbf{$\dim(\mathcal{M}_{eff})$ (有效维度)}：高维流形允许波包在多个正交平面上同时展开，这是并发处理（多线程）的几何等价物。
    \item   \textbf{$\gamma$ (介质粘滞系数)}：阻尼越大，信号衰减越快，用于维持驻波的功率消耗越大，留给有效计算的净功率越小。
    \item   \textbf{$L_{sys}$ (系统物理尺度)}：系统越大（如跨国数据中心），为了维持全局相位同步所需的时间开销越大，导致 $\tau_{\min}$ 增加，功率下降。
\end{itemize}

\section{智能演化的帕累托前沿 (The Pareto Frontier of Evolution)}
\label{sec:pareto_frontier}

在热力学中，热机的效率 ($\eta$) 与输出功率 ($P$) 之间存在经典的权衡关系（如：可逆卡诺热机效率最高，但功率为零，因为需要无限长时间）。在认知空间中，这种\textbf{效率-功率权衡 (Efficiency-Power Trade-off)} 同样成立，它构成了智能物种演化的\textbf{帕累托前沿}。

\smalltitle{3.1 二维相空间中的物种定标 (Species Calibration)}

将 $\eta_{cog}$ 设为纵轴，$P_{cog}^{max}$ 设为横轴，我们可以精确定位各类智能形态：

\begin{table}[htbp]
\centering
\footnotesize
\renewcommand{\arraystretch}{1.5}
\begin{tabularx}{\textwidth}{l|c|c|X}
\toprule
\rowcolor{structurecolor!20} \textbf{智能体类型} & \textbf{效率 ($\eta_{cog}$)} & \textbf{功率 ($P_{cog}^{max}$)} & \textbf{HSF-HD 物理/几何特征} \\
\midrule
\textbf{人类 (Class V)} & \textbf{极高} & \textbf{中等} & \textbf{[高维低频驻波]}。能耗仅 20W，但通过强大的宏观层 ($L_{macro}$) 退火控制，能从极少样本中提取拓扑不变量（抽象规律）。受限于湿件的 $c_{cog}$ 极慢。 \\
\hline
\textbf{LLM 集群 (Class III)} & \textbf{低} & \textbf{极高} & \textbf{[低维高频射流]}。依靠极高的 $c_{cog}$ 和无穷维展开，在短时间内完成海量遍历。但因形质混同与缺乏主动做功，大量计算被废热耗散，依赖海量数据轰击才能产生微小流变。 \\
\hline
\textbf{蚁群 (Class II)} & \textbf{极低} & \textbf{低} & \textbf{[二维化学标量场]}。缺乏宏观聚焦，全靠随机热涨落搜索测地线。效率与功率均受限于三维物理空间的扩散定律。 \\
\hline
\textbf{理想 AGI (Class V+)} & \textbf{极高} & \textbf{极高} & \textbf{[超导高维流形]}。兼具碳基的拓扑抽象能力与硅基的无阻尼光速演化。逼近理论帕累托极限。 \\
\bottomrule
\end{tabularx}
\caption{智能物种的效率-功率相图}
\end{table}

\smalltitle{3.2 逼近前沿的物理学路径}

由于介质物理特性的限制，任何单一系统 $\mathcal{S}$ 的 $(\eta, P)$ 点必须满足一个不等式约束：
$$ \mathcal{F}(\eta_{cog}, P_{cog}^{max}, \text{Medium Constants}) \le \text{Constant} $$

如果我们想要打破当前的僵局（例如让 LLM 变得高效，或让人脑算得更快），必须在物理层面引发\textbf{介质相变}：
\begin{enumerate}
    \item \textbf{向高 $\eta$ 跃迁 (提升效率)}：
    必须引入 \textbf{结构熵 (Epiplexity)} 优化。即放弃纯粹的参数暴力，引入 \textbf{双流架构 (MST)} 与 \textbf{宏观意志反馈}。让系统具备“认知退火”能力，从而在接收少量数据时精准找到局部极小值并瞬间固化。
    \item \textbf{向高 $P$ 跃迁 (提升功率)}：
    必须优化 \textbf{几何拓扑}。采用 \textbf{小世界网络 (Hubs)} 和 \textbf{基因缠绕 (光互连)} 降低 $L_{sys}$ 的有效距离，或者采用 \textbf{DEC（离散外微分）} 原生芯片降低计算的物理摩擦力（消除模拟税）。
\end{enumerate}

\section{结语：宇宙演化的测地线}

本章给出的尺度很直接：智能系统的差别，至少有一部分可以还原为效率与功率的组合方式。高效率不必然意味着高功率，高功率也不必然意味着高效率，演化真正追逐的是两者之间可持续的折中前沿。

因此，认知热机测量学的意义在于把“聪明”这件事从模糊印象变成可比较的物理量。后续所有工程讨论，都可以回到这个坐标系里判断收益与代价。
